\documentclass[11pt]{letter}

\usepackage[margin=0.75in]{geometry}
\usepackage{setspace}
\usepackage{hyperref}
\usepackage{xcolor}
\usepackage{enumitem}
\usepackage{parskip}

\hypersetup{
    colorlinks=true,
    urlcolor=black,
    linkcolor=black
}

\setlength{\parindent}{0pt}
\setstretch{1.05}
\date{}
\begin{document}

\begin{letter}{}

\opening{Dear Hiring Committee,}

I saw the job post on \textbf{LinkedIn} and thought I'd be a good fit for it.

I am a Master's student in Computer Science at the University of Cincinnati, graduating December 2026, looking for a role as an \textbf{AI Deep Learning/ML Engineer} at \textbf{Mentium Technologies}. With hands on PyTorch, TensorFlow, and OpenCV experience building and evaluating vision model pipelines, plus production software engineering experience, I am excited about the opportunity to contribute to Mentium's deep learning library for vision AI models on your in-memory compute hardware.

\textbf{Here are reasons that make me a great fit for the role:}

\begin{itemize}[leftmargin=*, itemsep=0.45em]
    \item \textbf{Vision model understanding:} As a Graduate Researcher at the University of Cincinnati, I built a PyTorch-based evaluation harness benchmarking a GPT-4o plus GroundingDINO (object detection) plus CLIP (Vision Transformer) pipeline against a proposed method across 53 runs on 20+ Android apps, achieving a 25.7x latency reduction while preserving task-completion fidelity, work that required close attention to inference time and model efficiency relevant to real-time edge deployment.

    \item \textbf{Python and numerical computing proficiency:} At Siemens Industry Software, I built modular Python automation (Boto3, object-oriented design) that replaced legacy scripts and cut deployment time by 40\%, and I regularly use NumPy and pandas across my ML projects for data processing and model evaluation.

    \item \textbf{Modular development and collaborative coding:} I migrated 40+ production environments to Kubernetes using ArgoCD, Helm, and GitLab CI/CD, collaborating through GitHub and producing clear technical documentation, and served as Quality Coach and SonarQube Specialist, raising automated code coverage from 35\% to 78\%.

    \item \textbf{Hardware and performance awareness:} My research work involved measuring inference time and latency tradeoffs directly, and my NYC Mobility Forecasting project applied PyTorch and scikit-learn to a 2.3M-row dataset, cutting RMSE 43\% over a baseline model, experience that translates well to reasoning about model size and precision tradeoffs for edge deployment.
\end{itemize}

I am available to work with you after my graduation in December 2026. As a Master's student at the University of Cincinnati and someone who values professionalism and growth, I am enthusiastic about contributing to Mentium's mission of advancing vision AI on ultra-low power hardware.

Please find my CV attached below. I look forward to hearing from you. I would be excited to help build a client-facing deep learning library that brings real-time vision AI to the edge.

{\setlength{\parskip}{0pt}
Regards,\\
Tanmay Bhuskute,\\
bhuskutv@mail.uc.edu $\mid$ tanmay.v.bhuskute@gmail.com $\mid$ +1 (732) 858-3027,\\
3424 Brookline Ave, Flat 19, Cincinnati, OH 45220.
}

\end{letter}

\end{document}
